% Généré par tools/build.py — ne pas éditer à la main
\section{Problèmes de synthèse, raisonnement et automatismes}

\fichehead{mar.~27~avr.~2027}{Problèmes de synthèse croisant plusieurs champs du programme}
\begin{fiche}
\qitem{Échauffement · Flash}{EQ}{Résoudre $5x-3=2x+9$.}{$3x=12$, $x=4$.}{}
\qitem{Cœur · Problème}{PRO}{Un prix subit une hausse de $20\,\%$ puis une remise fixe de $1\,000$ Ar et vaut alors $11\,000$ Ar. Quel était le prix initial\,?}{$1{,}2p-1\,000=11\,000$, $p=\dfrac{12\,000}{1{,}2}=10\,000$ Ar.}{}
\qitem{Cœur + · Méthode}{REG}{$f(x)=2x+1$ et $g(x)=-x+7$. Trouver le point d'intersection des courbes puis résoudre $f(x)<g(x)$.}{$2x+1=-x+7$~: $x=2$, point $(2;5)$\,; $f(x)<g(x)\Leftrightarrow3x<6\Leftrightarrow x<2$.}{}
\qitem{Retour · Flash}{QP}{Une recette pour $4$ personnes demande $250$ g de farine. Quelle quantité pour $6$ personnes\,?}{$250\div4\times6=375$ g.}{}
\qitem{Fil rouge · Flash}{GEO}{Une figure d'aire $4$ cm$^2$ est agrandie avec un rapport $3$. Quelle est l'aire de l'agrandissement\,?}{Aire $\times3^2$~: $36$ cm$^2$.}{Piège~: multiplier l'aire par $3$ au lieu de $3^2$.}
\end{fiche}

\fichehead{jeu.~29~avr.~2027}{Raisonnement par l'absurde et par disjonction de cas}
\begin{fiche}
\qitem{Échauffement · Flash}{DIV}{Quelle est la parité de $n(n+1)$ pour un entier $n$\,?}{Toujours pair~: l'un des deux entiers consécutifs est pair.}{}
\qitem{Cœur · Disjonction de cas}{DEM}{Montrer que $n(n+1)$ est pair pour tout entier $n$.}{Si $n$ est pair, le produit contient un facteur pair. Si $n$ est impair, $n+1$ est pair. Dans les deux cas, le produit est pair.}{}
\qitem{Cœur + · Absurde}{DEM}{Montrer par l'absurde que la somme d'un rationnel $r$ et d'un irrationnel $a$ est irrationnelle.}{Supposons $r+a=q\in\Q$\,; alors $a=q-r\in\Q$ (différence de deux rationnels)~: contradiction avec $a$ irrationnel.}{}
\qitem{Retour · Formule}{IMG}{$g(x)=2x+5$. Image de $-1$\,? Antécédent de $11$\,?}{$g(-1)=3$\,; $2x+5=11\Leftrightarrow x=3$.}{}
\qitem{Fil rouge · Formule → graphique}{REG}{Placer les points d'abscisses $-2\,;\,-1\,;\,0\,;\,1\,;\,2$ de la courbe de $f(x)=x^2$\,; donner leurs ordonnées.}{$4\ ;\ 1\ ;\ 0\ ;\ 1\ ;\ 4$~: les points sont symétriques par rapport à l'axe des ordonnées.}{}
\end{fiche}

\fichehead{lun.~3~mai~2027}{Connecteurs logiques, contraposée et contre-exemple en situation}
\begin{fiche}
\qitem{Échauffement · Flash}{DEM}{Écrire la négation de «~$x>2$~» puis de «~$x>2$ ou $x<-1$~».}{$x\leqslant2$\,; $x\leqslant2$ et $x\geqslant-1$, soit $-1\leqslant x\leqslant2$.}{}
\qitem{Cœur · Méthode}{DEM}{Écrire la contraposée de «~si $n^2$ est pair, alors $n$ est pair~».}{«~Si $n$ est impair, alors $n^2$ est impair~» (équivalente à l'énoncé initial).}{}
\qitem{Cœur + · Contre-exemple}{DEM}{Réfuter~: «~tout nombre premier est impair~».}{Contre-exemple~: $2$ est premier et pair.}{}
\qitem{Retour · Vrai/faux}{PNP}{Vrai ou faux\,? Justifier. Le périmètre d'un carré est proportionnel à la longueur de son côté.}{Vrai. $P=4c$~: coefficient de proportionnalité $4$.}{}
\qitem{Fil rouge · Lecture de code}{ALG}{Que renvoie $h(3)$\,? Écrire $h(x)$ sous forme d'une expression.}{$h(3)=5\times5=25$\,; $h(x)=(x+2)^2$.}{}
\end{fiche}
